\chapter{Matrices}\label{ch:b1:matrices}

Een matrix is een \omterm{def:b1:linmaps:def}{lineaire afbeelding}, opgeschreven in \omterm{prop:b1:vspaces:coordinates}{coördinaten}.
Dit hoofdstuk zet het woordenboek op --- samenstelling wordt
matrixproduct, bijectiviteit wordt inverteerbaarheid,
\omterm{def:b1:matrices:changeofbasis}{basisverandering} wordt conjugatie --- en de algoritmische kant:
\omterm{met:b1:matrices:gauss}{rij-operaties}, berekening van rangen en inversen. Matrices, voor het
eerst ontmoet in het bovenbouwvolume, worden nu gefundeerd in de
theorie van
\cref{ch:b1:vspaces,ch:b1:findim,ch:b1:linmaps}.

\section{Matrices en lineaire afbeeldingen}

\begin{definition}\label{def:b1:matrices:def}
$\mathcal{M}_{n,p}(K)$ is de \omterm{def:b1:vspaces:def}{vectorruimte} van de $n \times p$
tabellen $A = (a_{ij})$ scalairen ($i$: rij, $j$: kolom), van
dimensie $np$ (\omterm{def:b1:vspaces:free}{basis}: de matrices $E_{ij}$ met één enkele $1$).
Gegeven \omterm{def:b1:vspaces:free}{bases} $\mathcal{B} = (e_1, \dots, e_p)$ van $E$ en
$\mathcal{C}$ van $F$ ($\dim F = n$), is de \emph{matrix van $u \in
\mathcal{L}(E, F)$} de tabel waarvan de $j$-de kolom de \omterm{prop:b1:vspaces:coordinates}{coördinaten}
van $u(e_j)$ in $\mathcal{C}$ opsomt:
\[
\operatorname{Mat}_{\mathcal{B},\mathcal{C}}(u) = (a_{ij}),
\qquad u(e_j) = \sum_{i=1}^{n} a_{ij}\, f_i .
\]
De \omterm{def:b1:logic:map}{afbeelding} $u \mapsto
\operatorname{Mat}_{\mathcal{B},\mathcal{C}}(u)$ is een isomorfisme
van $\mathcal{L}(E, F)$ op $\mathcal{M}_{n,p}(K)$
(\cref{prop:b1:linmaps:basis}: een \omterm{def:b1:linmaps:def}{lineaire afbeelding} is precies
een keuze van \omterm{def:b1:linmaps:kerim}{beelden} van de $e_j$).
\end{definition}

\begin{example}[De afgeleide, als matrix]
\label{ex:b1:matrices:derivative}
Zij $D(P) = P'$ op $\R_3[X]$. In de monomiale \omterm{def:b1:vspaces:free}{basis} $(1, X, X^2,
X^3)$: $D(1) = 0$, $D(X) = 1$, $D(X^2) = 2X$, $D(X^3) = 3X^2$, dus
\[
\operatorname{Mat}(D) =
\begin{pmatrix}
0 & 1 & 0 & 0\\
0 & 0 & 2 & 0\\
0 & 0 & 0 & 3\\
0 & 0 & 0 & 0
\end{pmatrix}.
\]
In de \emph{gedeelde} \omterm{def:b1:vspaces:free}{basis} $\bigl(1,\ X,\ \frac{X^2}2,\
\frac{X^3}6\bigr)$ wordt elke basisvector op de vorige afgebeeld
($D\bigl(\frac{X^k}{k!}\bigr) = \frac{X^{k-1}}{(k-1)!}$), en wordt
de matrix de zuivere verschuiving: enen op de superdiagonaal,
nullen elders. Twee moralen: de matrix hoort bij het \emph{paar}
(\omterm{def:b1:logic:map}{afbeelding}, \omterm{def:b1:vspaces:free}{basis}), niet bij de \omterm{def:b1:logic:map}{afbeelding} alleen; en een goede
\omterm{def:b1:vspaces:free}{basis} maakt structuur in één oogopslag zichtbaar --- de
verschuivingsvorm toont ogenblikkelijk dat $D^4 = 0$ op $\R_3[X]$,
waarbij elke macht van de matrix haar diagonaal enen één stap
verder naar buiten duwt.
\end{example}

\begin{definition}[Product]\label{def:b1:matrices:product}
Voor $A \in \mathcal{M}_{n,p}$ en $B \in \mathcal{M}_{p,q}$:
\[
(AB)_{ik} = \sum_{j=1}^{p} a_{ij}\, b_{jk}
\qquad (1 \leq i \leq n,\ 1 \leq k \leq q).
\]
Dit is precies de matrix van de samenstelling:
$\operatorname{Mat}(v \circ u) = \operatorname{Mat}(v)\,
\operatorname{Mat}(u)$ (met in het midden overeenstemmende \omterm{def:b1:vspaces:free}{bases}).
Evenzo is, als $X$ de kolom \omterm{prop:b1:vspaces:coordinates}{coördinaten} van $x$ is, de kolom van
$u(x)$ gelijk aan $AX$.
\end{definition}

\begin{proof}[Bewijs van de formule voor de samenstelling]
\[
v(u(e_k)) = v\Bigl(\sum_j b_{jk} f_j\Bigr) = \sum_j b_{jk}\, v(f_j)
= \sum_j b_{jk} \sum_i a_{ij}\, g_i
= \sum_i \Bigl(\sum_j a_{ij} b_{jk}\Bigr) g_i . \qedhere
\]
\end{proof}

\begin{proposition}[De algebra $\mathcal{M}_n(K)$]\label{prop:b1:matrices:ring}
De vierkante matrices $\mathcal{M}_n(K)$ vormen een (voor $n \geq
2$ niet-commutatieve) \omterm{def:b1:structures:ring}{ring}, met eenheidselement $I_n$; haar \omterm{def:b1:structures:group}{groep}
van eenheden is de \emph{algemene lineaire groep}\index{algemene
lineaire groep} $GL_n(K)$, die overeenkomt met de \omterm{def:b1:logic:inj}{bijectieve}
endomorfismen. Voor $A, B \in \mathcal{M}_n(K)$ geldt
\[
AB = I_n \implies A \in GL_n(K) \text{ en } B = A^{-1}
\]
(eenzijdige inversen zijn tweezijdig, volgens
\cref{cor:b1:linmaps:samedim}).
\end{proposition}

\begin{proof}
De ringaxioma's reizen mee vanuit $\mathcal{L}(E)$ langs het
isomorfisme van \cref{def:b1:matrices:def}: dat zet samenstelling om
in product en som in som, zodat de associativiteit, de
distributiviteit en de rol van $I_n$ worden geërfd van de
overeenkomstige feiten over \omterm{def:b1:logic:map}{afbeeldingen}, zonder enige verificatie
per ingang. Niet-commutativiteit: $E_{12}E_{21} = E_{11} \neq E_{22}
= E_{21}E_{12}$. Is $AB = I_n$, dan voldoet het endomorfisme $a$ van
$A$ aan $a \circ b = \mathrm{id}$, zodat $a$ \omterm{def:b1:logic:inj}{surjectief} is ($x =
a(b(x))$ toont een \omterm{def:b1:logic:map}{origineel} van elke $x$), en dus \omterm{def:b1:logic:inj}{bijectief} in
\omterm{def:b1:findim:def}{eindige dimensie} (\cref{cor:b1:linmaps:samedim}); links met $a^{-1}$
samenstellen van $a \circ b = \mathrm{id}$ geeft $b = a^{-1}$, en
dan ook $b\circ a = \mathrm{id}$: de eenzijdige inverse was al die
tijd tweezijdig --- een strikt \omterm{def:b1:findim:def}{eindigdimensionale} gunst.
\end{proof}

\begin{definition}[Getransponeerde; spoor]\label{def:b1:matrices:transpose}
De \emph{getransponeerde} van $A = (a_{ij}) \in \mathcal{M}_{n,p}$
is $A^{\mathsf T} = (a_{ji}) \in \mathcal{M}_{p,n}$; zij voldoet aan
$(AB)^{\mathsf T} = B^{\mathsf T} A^{\mathsf T}$ en $(A^{\mathsf
T})^{\mathsf T} = A$. Het \emph{spoor}\index{spoor} van een
vierkante matrix is $\operatorname{tr} A = \sum_i a_{ii}$; het is
\omterm{def:b1:linmaps:def}{lineair}, en
\[
\operatorname{tr}(AB) = \operatorname{tr}(BA)
\qquad (A \in \mathcal{M}_{n,p},\ B \in \mathcal{M}_{p,n}).
\]
\end{definition}

\begin{proof}[Bewijs van de identiteit voor het spoor]
$\operatorname{tr}(AB) = \sum_i \sum_j a_{ij} b_{ji}$ en
$\operatorname{tr}(BA) = \sum_j \sum_i b_{ji} a_{ij}$: dezelfde
dubbele som.
\end{proof}

\begin{example}[Het spoor aan het werk]
\label{ex:b1:matrices:tracework}
De \omterm{def:b1:linmaps:projection}{projectie} uit \cref{ch:b1:linmaps} op
$\operatorname{Vect}(1,1)$ langs $\operatorname{Vect}(0,1)$, $p(x,
y) = (x, x)$, heeft in de canonieke \omterm{def:b1:vspaces:free}{basis} de matrix $A =
\begin{pmatrix} 1 & 0\\ 1 & 0\end{pmatrix}$: inderdaad is $A^2 =
A$, en
\[
\operatorname{tr} A = 1 = \operatorname{rk} A ,
\]
wat \cref{exo:b1:matrices:8} illustreert: voor idempotenten
\emph{telt} het \omterm{def:b1:matrices:transpose}{spoor} de dimensie van het \omterm{def:b1:linmaps:kerim}{beeld}, in welke scheve
\omterm{def:b1:vspaces:free}{basis} de matrix ook is geschreven. Het mechanisme van de
invariantie is de identiteit $\operatorname{tr}(AB) =
\operatorname{tr}(BA)$:
\[
\operatorname{tr}\bigl(P^{-1}(AP)\bigr) =
\operatorname{tr}\bigl((AP)P^{-1}\bigr) = \operatorname{tr} A ,
\]
zodat alle matrices die gelijkvormig zijn met $A$ haar \omterm{def:b1:matrices:transpose}{spoor}
delen --- de eerste \emph{numerieke invariant} van een
endomorfisme, waar in \cref{ch:b1:det} de determinant zich bij zal
voegen (het paar $(s, p)$ van de weekendopgave hieronder).
\end{example}

\begin{example}[Symmetrisch plus antisymmetrisch]
\label{ex:b1:matrices:symsplit}
Noem $A$ \emph{symmetrisch} wanneer $A^{\mathsf T} = A$, en
\emph{antisymmetrisch} wanneer $A^{\mathsf T} = -A$. Elke vierkante
matrix splitst zich eenduidig als de ene plus de andere:
\[
A = \underbrace{\frac{A + A^{\mathsf T}}{2}}_{\text{symmetrisch}}
+ \underbrace{\frac{A - A^{\mathsf T}}{2}}_{\text{antisymmetrisch}},
\]
en een matrix die beide is, is nul ($A = -A$): de twee
\omterm{def:b1:logic:sets}{verzamelingen} zijn \omterm{def:b1:vspaces:sum}{complementaire deelruimten} van
$\mathcal{M}_n(K)$ --- het exacte analogon van de splitsing in even
en oneven functies (\cref{ex:b1:vspaces:evenodd}), waarbij het
transponeren de rol van $x \mapsto -x$ speelt. Dimensies: een
symmetrische matrix is \omterm{def:b1:vspaces:free}{vrij} op en boven de diagonaal, een
antisymmetrische strikt erboven (met nuldiagonaal):
\[
\frac{n(n+1)}{2} + \frac{n(n-1)}{2} = n^2 ,
\]
en dat de telling klopt, is de bevestiging van de directheid door
Grassmann. Voor $n = 2$: $\begin{pmatrix} 1 & 5\\ 1 &
2\end{pmatrix} = \begin{pmatrix} 1 & 3\\ 3 & 2\end{pmatrix} +
\begin{pmatrix} 0 & 2\\ -2 & 0\end{pmatrix}$. Symmetrische matrices
keren terug als de gegevens van de tweede \omterm{def:b1:derivative:def}{afgeleide} in
\cref{ch:b1:multivar} (het drietal $r, s, t$ van Monge), en de
symmetrisch-orthogonale worden in \cref{exo:b1:euclid:12}
geclassificeerd.
\end{example}

\section{Basisverandering}

\begin{definition}\label{def:b1:matrices:changeofbasis}
Zij $\mathcal{B}, \mathcal{B}'$ \omterm{def:b1:vspaces:free}{bases} van $E$. De matrix van
\emph{basisverandering}\index{basisverandering} $P =
P_{\mathcal{B}\to\mathcal{B}'}$ heeft als kolommen de \omterm{prop:b1:vspaces:coordinates}{coördinaten}
van de \emph{nieuwe} basisvectoren in de \emph{oude} \omterm{def:b1:vspaces:free}{basis}. Zij is
inverteerbaar, $P^{-1} = P_{\mathcal{B}'\to\mathcal{B}}$, en de
\omterm{prop:b1:vspaces:coordinates}{coördinaten} transformeren volgens $X = PX'$ (oud $=$ $P\,\cdot$
nieuw).
\end{definition}

\begin{example}[De matrix van basisverandering lezen]
\label{ex:b1:matrices:changerun}
In $\R^2$, van de canonieke $\mathcal B$ naar $\mathcal B' =
\bigl((1,1), (1,-1)\bigr)$:
\[
P = P_{\mathcal B\to\mathcal B'} =
\begin{pmatrix} 1 & 1\\ 1 & -1 \end{pmatrix}
\]
(de nieuwe vectoren geschreven in oude \omterm{prop:b1:vspaces:coordinates}{coördinaten}, kolom voor
kolom). De vector met oude \omterm{prop:b1:vspaces:coordinates}{coördinaten} $X = (3, 1)^{\mathsf T}$
heeft nieuwe \omterm{prop:b1:vspaces:coordinates}{coördinaten} $X' = P^{-1}X = \frac12(3 + 1,\ 3 -
1)^{\mathsf T} = (2, 1)^{\mathsf T}$: inderdaad is $2(1,1) +
1(1,-1) = (3,1)$. Let op de richting --- de matrix $P$ wordt uit de
\emph{nieuwe} \omterm{def:b1:vspaces:free}{basis} gebouwd maar zet \emph{nieuw om in oud} ($X =
PX'$); van oud naar nieuw gaan kost de inverse. De controle
$2(1,1) + (1,-1) = (3,1)$ opschrijven na elke omzetting vangt de
fout met de omgekeerde $P$ op, die de meest voorkomende fout van
het hoofdstuk is.
\end{example}

\begin{theorem}[Basisverandering voor een afbeelding]\label{thm:b1:matrices:conjugation}
Zij $u \in \mathcal{L}(E)$ met matrix $A$ in $\mathcal{B}$ en $A'$
in $\mathcal{B}'$, en $P = P_{\mathcal{B}\to\mathcal{B}'}$. Dan is
\[
A' = P^{-1} A\, P .
\]
Twee matrices die zo verwant zijn, heten
\emph{gelijkvormig}\index{gelijkvormige matrices}. (Voor $u \colon E
\to F$ met twee paren \omterm{def:b1:vspaces:free}{bases} luidt de formule $A' = Q^{-1} A P$ ---
\emph{equivalente} matrices\index{equivalente matrices}.)
\end{theorem}

\begin{proof}
Voor elke $x$ geldt $X = PX'$ en voldoet het \omterm{def:b1:linmaps:kerim}{beeld} aan $Y = AX$, $Y
= PY'$. Dus $PY' = APX'$, dat wil zeggen $Y' = (P^{-1}AP)X'$ voor
alle $X'$: de matrix van $u$ in de nieuwe \omterm{def:b1:vspaces:free}{basis} is $P^{-1}AP$ (neem
voor $X'$ de canonieke kolommen).
\end{proof}

\begin{example}[Een goede basis maakt een afbeelding doorzichtig]
\label{ex:b1:matrices:conjugationrun}
Zij $u(x, y) = (y, x)$ (verwisseling), met matrix $A =
\begin{pmatrix} 0 & 1\\ 1 & 0\end{pmatrix}$ in de canonieke \omterm{def:b1:vspaces:free}{basis}.
In de \omterm{def:b1:vspaces:free}{basis} $\mathcal B' = \bigl((1,1), (1,-1)\bigr)$:
\[
P = \begin{pmatrix} 1 & 1\\ 1 & -1 \end{pmatrix},
\qquad
P^{-1} = \frac12\begin{pmatrix} 1 & 1\\ 1 & -1 \end{pmatrix},
\qquad
P^{-1} A P = \begin{pmatrix} 1 & 0\\ 0 & -1 \end{pmatrix}.
\]
Er was eigenlijk geen matrixproduct nodig: $u$ houdt $(1,1)$ vast
en keert $(1,-1)$ om, dus \emph{moet} haar matrix in $\mathcal B'$
gelijk zijn aan $\operatorname{diag}(1, -1)$ --- de verwisseling is
de spiegeling in de rechte $y = x$. Voor een gegeven endomorfisme
een \omterm{def:b1:vspaces:free}{basis} vinden waarin zijn matrix diagonaal wordt, is het
centrale probleem van het volume van bachelorjaar 2
(reductietheorie); de weekendopgave hieronder toont hoe ver
veeltermidentiteiten alleen al reiken.
\end{example}

\begin{example}[Basisverandering, achterstevoren gelopen]
\label{ex:b1:matrices:reverserun}
De \omterm{def:b1:linmaps:projection}{projectie} op $F = \operatorname{Vect}(1,1)$ langs $G =
\operatorname{Vect}(1,-1)$ heeft, in de aangepaste \omterm{def:b1:vspaces:free}{basis} $\mathcal
B' = \bigl((1,1),(1,-1)\bigr)$, de doorzichtige matrix $A' =
\operatorname{diag}(1, 0)$. Om haar matrix in de canonieke \omterm{def:b1:vspaces:free}{basis} te
krijgen, laat men \cref{thm:b1:matrices:conjugation} achterstevoren
lopen, $A = P A' P^{-1}$:
\[
P = \begin{pmatrix} 1 & 1\\ 1 & -1\end{pmatrix},
\quad
P^{-1} = \frac12\begin{pmatrix} 1 & 1\\ 1 & -1\end{pmatrix},
\quad
A = P\begin{pmatrix} 1 & 0\\ 0 & 0\end{pmatrix}P^{-1}
= \frac12\begin{pmatrix} 1 & 1\\ 1 & 1\end{pmatrix}.
\]
Controle: $A^2 = A$ (idempotent), $\operatorname{tr} A = 1 =
\operatorname{rk} A$, en $A\binom{1}{1} = \binom11$,
$A\binom{1}{-1} = 0$, zoals voorgeschreven. Deze omgekeerde
richting --- ontwerp de matrix in de goede \omterm{def:b1:vspaces:free}{basis} en conjugeer terug
--- is hoe matrices van rotaties, spiegelingen en \omterm{def:b1:linmaps:projection}{projecties} in de
praktijk werkelijk worden gemaakt.
\end{example}

\begin{theorem}[Normaalvorm van de rang]\label{thm:b1:matrices:rank}
De \emph{rang} van een matrix (de rang van haar kolommen, equivalent
die van de bijbehorende \omterm{def:b1:linmaps:def}{lineaire afbeelding}) is de enige invariant
van de equivalentie: elke $A \in \mathcal{M}_{n,p}$ van rang $r$ is
equivalent met
\[
J_r = \begin{pmatrix} I_r & 0 \\ 0 & 0 \end{pmatrix},
\]
en $\operatorname{rk}(A^{\mathsf T}) = \operatorname{rk}(A)$: de
rijrang is gelijk aan de kolomrang.
\end{theorem}

\begin{proof}
Zij $u \colon E \to F$ van rang $r$. Kies een \omterm{def:b1:vspaces:sum}{complementaire} $S$ van
$\ker u$ ($\dim S = r$, \cref{thm:b1:linmaps:ranknullity}) met \omterm{def:b1:vspaces:free}{basis}
$(e_1, \dots, e_r)$, aangevuld met een \omterm{def:b1:vspaces:free}{basis} van $\ker u$ tot een
\omterm{def:b1:vspaces:free}{basis} van $E$; de \omterm{def:b1:linmaps:kerim}{beelden} $f_i = u(e_i)$, $i \leq r$, vormen een
\omterm{def:b1:vspaces:free}{basis} van $\operatorname{im} u$ (de beperking is een isomorfisme),
aangevuld tot een \omterm{def:b1:vspaces:free}{basis} van $F$. In die \omterm{def:b1:vspaces:free}{bases} is de matrix van $u$
precies $J_r$. Dus $A = Q J_r P^{-1}$ met inverteerbare $P, Q$.

Transponeren: $A^{\mathsf T} = (P^{-1})^{\mathsf T} J_r^{\mathsf T}
Q^{\mathsf T}$ met $J_r^{\mathsf T}$ van dezelfde gedaante (rang
$r$) en met inverteerbare buitenste factoren (de getransponeerde van
een inverteerbare is inverteerbaar, uit $(AB)^{\mathsf T} =
B^{\mathsf T}A^{\mathsf T}$ toegepast op $AA^{-1} = I$):
$\operatorname{rk} A^{\mathsf T} = r$.
\end{proof}

\section{Rij-operaties}

\begin{method}[Gauss-eliminatie op matrices]\label{met:b1:matrices:gauss}
De drie \emph{elementaire rij-operaties}\index{rij-operaties} ---
twee rijen verwisselen, een rij met $\lambda \neq 0$
vermenigvuldigen, een veelvoud van een rij bij een andere optellen
--- veranderen de rang niet (elk is een linkse vermenigvuldiging met
een inverteerbare matrix). Algoritme: maak een spil (de meest linkse
ingang ongelijk aan nul), veeg haar kolom eronder schoon, en ga naar
de volgende rij en kolom; het aantal spillen van de resulterende
trapvorm is de rang.

\emph{Berekening van de inverse:} laat het algoritme lopen op het
blok $(A \mid I_n)$ tot het linkerblok $I_n$ wordt (mogelijk dan en
slechts dan als $A$ inverteerbaar is); het rechterblok is dan
$A^{-1}$ --- immers, het product van de gebruikte elementaire
matrices is gelijk aan $A^{-1}$.
\end{method}

\begin{example}\label{ex:b1:matrices:inverse}
$A = \begin{pmatrix} 1 & 2 \\ 3 & 4 \end{pmatrix}$: reduceer $(A
\mid I_2)$:
\[
\begin{pmatrix} 1 & 2 & 1 & 0\\ 3 & 4 & 0 & 1 \end{pmatrix}
\to
\begin{pmatrix} 1 & 2 & 1 & 0\\ 0 & -2 & -3 & 1 \end{pmatrix}
\to
\begin{pmatrix} 1 & 0 & -2 & 1\\ 0 & 1 & \tfrac32 & -\tfrac12
\end{pmatrix},
\]
(bewerkingen: $L_2 \leftarrow L_2 - 3L_1$; daarna $L_1 \leftarrow
L_1 + L_2$, $L_2 \leftarrow -\frac12 L_2$). Dus $A^{-1} =
\begin{pmatrix} -2 & 1 \\ \tfrac32 & -\tfrac12\end{pmatrix}$.
\emph{Controle:} $AA^{-1} = I_2$.
\end{example}

\begin{example}[Rang met een parameter, alleen met rijen]
\label{ex:b1:matrices:rankparameter}
Voor $m \in \R$ de rang van $M_m = \begin{pmatrix} 1 & 1 & m\\ 1 &
m & 1\\ m & 1 & 1\end{pmatrix}$. Reduceer: $L_2 \leftarrow L_2 -
L_1$ en $L_3 \leftarrow L_3 - mL_1$ geven de rijen
\[
(1,\ 1,\ m), \qquad (0,\ m - 1,\ 1 - m), \qquad
(0,\ 1 - m,\ 1 - m^2).
\]
\emph{Geval $m = 1$}: de laatste twee rijen worden nul --- één
spil, $\operatorname{rk} M_1 = 1$ (alle drie de oorspronkelijke
rijen waren gelijk). \emph{Geval $m \neq 1$}: schaal $L_2$ met
$\frac1{m-1}$ en $L_3$ met $\frac1{1-m}$ om $(0, 1, -1)$ en $(0,
1, 1 + m)$ te krijgen, en daarna $L_3 \leftarrow L_3 - L_2 = (0,
0, m + 2)$. Is $m = -2$: twee spillen, rang $2$; anders drie
spillen, rang $3$. Samengevat:
\[
\operatorname{rk} M_m =
\begin{cases}
1 & m = 1,\\
2 & m = -2,\\
3 & \text{anders}.
\end{cases}
\]
Dezelfde drempels zullen uit één determinantberekening in
\cref{ch:b1:det} rollen (de \omterm{def:b1:poly:def}{veelterm} $-(m+2)(m-1)^2$ van
\cref{exo:b1:det:7}) --- maar merk op wat de eliminatie geeft en de
determinant niet: de \emph{waarde} van de rang in de ontaarde
gevallen, en niet alleen het feit dat zij daalde.
\end{example}

\begin{example}[Machten berekenen]\label{ex:b1:matrices:powers}
$A = \begin{pmatrix} 1 & 1 \\ 0 & 1\end{pmatrix} = I + N$ met $N =
E_{12}$ en $N^2 = 0$. Omdat $I$ en $N$ commuteren, breekt het
binomium (\cref{prop:b1:structures:binomial}) af:
\[
A^k = I + kN = \begin{pmatrix} 1 & k \\ 0 & 1 \end{pmatrix}
\qquad (k \in \N, \text{ en } k \in \Z \text{ met } A^{-1} = I -
N).
\]
\end{example}

\begin{method}[$A^n$ berekenen: de drie wegen]
\label{met:b1:matrices:powers}
\begin{enumerate}
  \item \emph{Binomiale weg}: is $A = \lambda I + N$ met $N$
        nilpotent, dan breekt het binomium af
        (\cref{ex:b1:matrices:powers}, \cref{exo:b1:matrices:5});
        het is toepasbaar omdat $\lambda I$ met alles commuteert.
  \item \emph{Veeltermweg}: zoek een veeltermidentiteit waaraan
        $A$ voldoet (in dimensie $2$ altijd $A^2 = sA - pI$) en
        reduceer $X^n$ modulo die identiteit; de weekendopgave
        hieronder bouwt deze weg volledig op.
  \item \emph{Gelijkvormigheidsweg}: zoek een inverteerbare $P$ met
        $P^{-1}AP = D$ eenvoudig (diagonaal, of een verschuiving),
        bereken $D^n$, en draai terug: $A^n = P D^n P^{-1}$
        (\cref{thm:b1:matrices:conjugation},
        \cref{ex:b1:matrices:conjugationrun}); het systematisch
        zoeken naar zo'n $P$ is de reductietheorie van bachelorjaar
        2.
\end{enumerate}
Welke weg men ook neemt, controleer het resultaat bij $n = 0, 1,
2$: drie goedkope tests die vrijwel elke misstap opvangen.
\end{method}

\begin{remark}[Veelgemaakte fouten: de prijs van de
niet-commutativiteit]
\label{rem:b1:matrices:pitfalls}
Elke identiteit uit de scalaire algebra waarvan het bewijs
factoren herordent, sterft in $\mathcal{M}_n(K)$ met $n \geq 2$.
\emph{Kwadraten}: $(A + B)^2 = A^2 + AB + BA + B^2$, en het midden
stort alleen in tot $2AB$ als $AB = BA$
(\cref{exo:b1:matrices:1}).
\emph{Machten van producten}: $(AB)^k$ is $ABAB\cdots$, niet
$A^kB^k$.
\emph{Nuldelers}: $E_{12}E_{12} = 0$ met $E_{12} \neq 0$;
bijgevolg \emph{geen wegstrepen}: $AB = AC$ impliceert $B = C$
alleen wanneer $A$ inverteerbaar is (vermenigvuldig met $A^{-1}$
--- aan de juiste kant).
\emph{\omterm{def:b1:matrices:transpose}{Sporen}}: $\operatorname{tr}(AB) = \operatorname{tr}(BA)$
geldt altijd, maar $\operatorname{tr}(AB) \neq
\operatorname{tr}A\operatorname{tr}B$ in het algemeen (neem $A = B
= I_2$: $2 \neq 4$), en $\operatorname{tr}(ABC) =
\operatorname{tr}(BCA)$ (cyclisch) terwijl
$\operatorname{tr}(ACB)$ kan verschillen.
\emph{Getransponeerden keren om}: $(AB)^{\mathsf T} = B^{\mathsf
T}A^{\mathsf T}$ --- de omkering vergeten is de meest voorkomende
fout in berekeningen met orthogonaliteit (\cref{ch:b1:euclid}).
Bij twijfel: toets elke beweerde identiteit op $E_{12}$ en
$E_{21}$; het kleinste niet-commuterende paar weerlegt de meeste
valse formules in één regel.
\end{remark}

\begin{remark}[Waar het woordenboek heen gaat]
\label{rem:b1:matrices:whereused}
Het matrixwoordenboek wordt op elke resterende bladzijde van dit
volume gebruikt: \cref{ch:b1:det} hangt aan elke vierkante matrix
één getal dat over de inverteerbaarheid beslist, en lost $AX = B$
systematisch op; \cref{ch:b1:euclid} zondert de matrices af die
lengten bewaren (orthogonale matrices); en in
\cref{ch:b1:multivar} is het gedrag van tweede orde van een
functie van twee veranderlijken een symmetrische $2 \times 2$
matrix. Het \omterm{def:b1:matrices:transpose}{spoor}, hierboven bijna terloops ingevoerd, wordt een
krachtige invariant:
\cref{exo:b1:matrices:6,exo:b1:matrices:8} geven een eerste
voorproef, en het volume van bachelorjaar 2 bouwt er de theorie
van de eigenwaarden op. De weekendopgave ontwikkelt het andere
werkpaard: veeltermidentiteiten waaraan een matrix voldoet, die de
berekening van $A^n$ omzetten in een \omterm{pb:b1:matrices:1}{lineaire recurrentie} met twee
termen.
\end{remark}

\begin{remark}[Vooruitzichten binnen boek 3]
\label{rem:b1:matrices:perspectives}
Drie families matrices die hier zijn ingevoerd, hebben verderop in
dit volume een afspraak. \emph{Symmetrische matrices}
(\cref{ex:b1:matrices:symsplit}) dragen de gegevens van de tweede
orde van functies van twee veranderlijken: de toets van Monge uit
\cref{ch:b1:multivar} is een \omterm{def:b1:logic:statement}{uitspraak} over het tekengedrag van
een symmetrische $2\times2$ matrix, en haar determinant $rt - s^2$
wordt met de machinerie van \cref{ch:b1:det} berekend.
\emph{Orthogonale matrices} ($A^{\mathsf T}A = I$) zijn de
isometrieën van \cref{ch:b1:euclid}, waar de getransponeerde
eindelijk haar meetkundige betekenis krijgt: zij is de
algebraïsche schaduw van het inproduct. \emph{Inverteerbare
matrices} ontmoeten hun praktische toets in \cref{ch:b1:det} ---
één getal, $\det A \neq 0$ --- waarmee de zoektocht wordt \omterm{def:b1:topology:closed}{gesloten}
die dit hoofdstuk met rijreductie begon. \omterm{def:b1:matrices:transpose}{Spoor} en determinant
reizen daarna verder als het invariante paar $(s, p)$ van de
weekendopgave, helemaal tot in de theorie van de eigenwaarden van
bachelorjaar 2.
\end{remark}

\section{Oefeningen}

\begin{exercise}[$\star$]\label{exo:b1:matrices:1}
Zij $A = \begin{pmatrix} 1 & 2 \\ 0 & 1 \end{pmatrix}$ en $B =
\begin{pmatrix} 0 & 1 \\ 1 & 0\end{pmatrix}$. Bereken $AB$, $BA$,
$A^2 - B^2$ en $(A+B)(A-B)$; leg uit waarom de laatste twee
verschillen.
\end{exercise}

\begin{exercise}[$\star$]\label{exo:b1:matrices:2}
Bereken de rang van
\[
M = \begin{pmatrix}
1 & 2 & 3\\
2 & 4 & 6\\
1 & 1 & 1
\end{pmatrix},
\qquad
N = \begin{pmatrix}
1 & 1 & 0 & 2\\
0 & 1 & 1 & 1\\
1 & 2 & 1 & 3
\end{pmatrix}.
\]
\end{exercise}

\begin{exercise}[$\star$]\label{exo:b1:matrices:3}
Inverteer met rijreductie $A = \begin{pmatrix} 1 & 0 & 1\\ 2 & 1 &
1\\ 1 & 1 & 1 \end{pmatrix}$, en controleer met één product.
\end{exercise}

\begin{exercise}[$\star$]\label{exo:b1:matrices:4}
Schrijf de matrix, in de canonieke \omterm{def:b1:vspaces:free}{basis} van $\R_2[X]$, van het
endomorfisme $u(P) = P(X + 1)$. Leg zonder berekening uit waarom
zij inverteerbaar is, en geef de matrix van $u^{-1}$.
\end{exercise}

\begin{exercise}[$\star\star$]\label{exo:b1:matrices:5}
Zij $A = \begin{pmatrix} 2 & 1 \\ 0 & 2\end{pmatrix}$. Schrijf $A =
2I + N$, bereken $N^2$, en leid met het binomium $A^k$ af voor alle
$k \in \N$.
\end{exercise}

\begin{exercise}[$\star\star$]\label{exo:b1:matrices:6}
Bewijs dat er geen matrices $A, B \in \mathcal{M}_n(K)$ (met $K =
\R$ of $\C$) bestaan met $AB - BA = I_n$. \emph{(Neem \omterm{def:b1:matrices:transpose}{sporen}.)}
\end{exercise}

\begin{exercise}[$\star\star$]\label{exo:b1:matrices:7}
Een matrix $A$ heet \emph{nilpotent} wanneer $A^m = 0$ voor zekere
$m$. Bewijs dat $I - A$ dan inverteerbaar is, met
\[
(I - A)^{-1} = I + A + A^2 + \dots + A^{m-1} .
\]
Toepassing: inverteer $\begin{pmatrix} 1 & 2 & 3\\ 0 & 1 & 2\\ 0 &
0 & 1\end{pmatrix}$.
\end{exercise}

\begin{exercise}[$\star\star$]\label{exo:b1:matrices:8}
Zij $A \in \mathcal{M}_n(\R)$ met $A^2 = A$ (idempotent). Bewijs
dat $\operatorname{tr} A = \operatorname{rk} A$. \emph{(Vat $A$ op
als een \omterm{def:b1:linmaps:projection}{projectie} en kies een aangepaste \omterm{def:b1:vspaces:free}{basis};
\cref{thm:b1:matrices:conjugation} zegt dat het \omterm{def:b1:matrices:transpose}{spoor} niet van de
\omterm{def:b1:vspaces:free}{basis} afhangt, want $\operatorname{tr}(P^{-1}MP) =
\operatorname{tr} M$.)}
\end{exercise}

\begin{exercise}[$\star\star\star$]\label{exo:b1:matrices:9}
Zij $J \in \mathcal{M}_n(\R)$ de matrix met overal enen. Bereken
$J^2$, en leid daaruit voor $a, b \in \R$ de voorwaarde voor de
inverteerbaarheid van $M = aI + bJ$ af, samen met $M^{-1}$
\emph{(zoek een inverse van dezelfde vorm $\alpha I + \beta J$)}.
\end{exercise}

\begin{exercise}[$\star\star\star$]\label{exo:b1:matrices:10}
(Ongelijkheden voor de rang) Bewijs voor $A, B \in
\mathcal{M}_n(K)$:
\[
\operatorname{rk}(A + B) \leq \operatorname{rk} A +
\operatorname{rk} B,
\qquad
\operatorname{rk}(AB) \geq \operatorname{rk} A + \operatorname{rk}
B - n .
\]
\emph{(Pas voor de tweede --- de ongelijkheid van Sylvester --- de
dimensiestelling toe op de beperking van de \omterm{def:b1:logic:map}{afbeelding} van $A$ tot
$\operatorname{im} B$.)}
\end{exercise}

\begin{exercise}[$\star\star$]\label{exo:b1:matrices:11}
Zij $D = \operatorname{diag}(d_1, \dots, d_n)$ met de $d_i$
\emph{paarsgewijs verschillend}.
\begin{enumerate}
  \item Bewijs dat een matrix $A$ met $D$ commuteert dan en slechts
        dan als $A$ diagonaal is. \emph{(Vergelijk de ingangen
        $(i,j)$ van $AD$ en $DA$.)}
  \item Leid het \emph{centrum} van $\mathcal{M}_n(K)$ af: de
        matrices die met \emph{elke} matrix commuteren, zijn precies
        de scalaire matrices $\lambda I_n$. \emph{(Toets tegen $D$,
        en daarna tegen de matrices $E_{ij}$.)}
\end{enumerate}
\end{exercise}

\begin{exercise}[$\star\star\star$]\label{exo:b1:matrices:12}
(Matrices van rang één) Zij $A \in \mathcal{M}_n(K)$, $A \neq 0$.
\begin{enumerate}
  \item Bewijs dat $\operatorname{rk} A = 1$ dan en slechts dan als
        $A = CL$ voor een kolom $C \in \mathcal{M}_{n,1}$ ongelijk
        aan nul en een rij $L \in \mathcal{M}_{1,n}$ ongelijk aan
        nul.
  \item Bewijs voor zo'n $A$ dat $A^2 = (\operatorname{tr} A)\,A$;
        leid af dat een matrix van rang één nilpotent is dan en
        slechts dan als haar \omterm{def:b1:matrices:transpose}{spoor} nul is.
  \item Bewijs dat als $\operatorname{tr} A \neq -1$, de matrix
        $I_n + A$ inverteerbaar is met
        \[
        (I_n + A)^{-1} = I_n - \frac{1}{1 + \operatorname{tr}
        A}\,A ,
        \]
        en dat $I_n + A$ \emph{niet} inverteerbaar is wanneer
        $\operatorname{tr} A = -1$. \emph{(Zoek een vector die door
        $I_n + A$ wordt gedood.)}
\end{enumerate}
\end{exercise}

\section{Opgave: machten van een matrix via veeltermdeling}

\begin{problem}\label{pb:b1:matrices:1}
$A^{100}$ ingang voor ingang berekenen is hopeloos; het berekenen
via een veeltermidentiteit waaraan $A$ voldoet, kost drie regels.
Deze opgave bouwt de methode vanaf nul op: de euclidische deling
van $X^n$, de identiteit $A^2 - sA + pI = 0$ waaraan elke $2 \times
2$ matrix voldoet (de stelling van Cayley--Hamilton in dimensie
$2$), en het woordenboek tussen machten van matrices en \omterm{pb:b1:matrices:1}{lineaire
recurrenties} --- met de fibonaccigetallen als doorlopend
voorbeeld.
\index{stelling van Cayley--Hamilton (dimensie 2)}
\index{Fibonacci-getallen}\index{lineaire recurrentie}

\medskip
\noindent\textbf{Deel I --- Het rekenen met resten.} Leg $s, p \in
K$ vast en stel $D = X^2 - sX + p$.
\begin{enumerate}
  \item Verantwoord dat er voor elke $n \in \N$ eenduidige $Q_n \in
        K[X]$ en $(a_n, b_n) \in K^2$ bestaan met
        \[
        X^n = Q_n\,D + a_n X + b_n ,
        \]
        en bereken $(a_0, b_0)$ en $(a_1, b_1)$.
  \item Stel, door met $X$ te vermenigvuldigen en opnieuw te delen,
        de recursies
        \[
        a_{n+1} = s\,a_n + b_n,
        \qquad
        b_{n+1} = -p\,a_n
        \]
        vast, en leid af dat $a_{n+2} = s\,a_{n+1} - p\,a_n$: de rij
        coëfficiënten gehoorzaamt aan de \omterm{pb:b1:matrices:1}{lineaire recurrentie} die
        bij $D$ hoort.
  \item Stel dat $D$ twee verschillende wortels $\lambda \neq \mu$
        heeft. Bewijs door de delingsidentiteit te evalueren dat
        \[
        a_n = \frac{\lambda^n - \mu^n}{\lambda - \mu},
        \qquad
        b_n = \frac{\lambda\mu^n - \mu\lambda^n}{\lambda - \mu} .
        \]
  \item Stel dat $D = (X - \lambda)^2$. Bewijs met de \omterm{def:b1:derivative:def}{afgeleide} van
        de delingsidentiteit dat $a_n = n\lambda^{n-1}$ en $b_n =
        (1 - n)\lambda^{n}$.
  \item Toon aan dat het invullen van een vaste matrix $M \in
        \mathcal{M}_k(K)$ in \omterm{def:b1:poly:def}{veeltermen} sommen en producten
        respecteert: $(PQ)(M) = P(M)\,Q(M)$. Leid af dat als $D(M)
        = 0$, dan
        \[
        M^n = a_n\,M + b_n\,I \qquad (n \in \N).
        \]
\end{enumerate}

\medskip
\noindent\textbf{Deel II --- Dimensie 2: \omterm{def:b1:matrices:transpose}{spoor}, determinantgetal,
Cayley--Hamilton.} Stel voor $A = \begin{pmatrix} a & b\\ c &
d\end{pmatrix}$ dat $s = a + d = \operatorname{tr} A$ en $p = ad -
bc$ (het getal dat \cref{ch:b1:det} de determinant zal noemen).
\begin{enumerate}[resume]
  \item Ga door rechtstreekse berekening de \emph{identiteit van
        Cayley--Hamilton in dimensie $2$} na:
        \[
        A^2 - s\,A + p\,I_2 = 0 .
        \]
  \item Bewijs door rechtstreeks uit te werken dat $p$
        multiplicatief is: met de voor de hand liggende notatie,
        $p(AB) = p(A)\,p(B)$. Toon daarna aan dat $A$ inverteerbaar
        is dan en slechts dan als $p \neq 0$, en dat dan
        \[
        A^{-1} = \frac1p\,\bigl(s\,I_2 - A\bigr).
        \]
  \item Zij $A = \begin{pmatrix} 1 & 1\\ 0 & 2\end{pmatrix}$.
        Bereken $s$, $p$ en de wortels van $D$, en leid een \omterm{def:b1:topology:closed}{gesloten}
        formule voor $A^n$ af; toets haar aan een rechtstreekse
        berekening van $A^2$.
  \item Zij $A = \begin{pmatrix} 3 & 1\\ -1 & 1\end{pmatrix}$. Toon
        aan dat $D$ een dubbele wortel heeft en bereken $A^n$;
        controleer bij $n = 2$.
  \item Zij $F = \begin{pmatrix} 1 & 1\\ 1 & 0\end{pmatrix}$ en
        definieer de fibonaccigetallen door $F_0 = 0$, $F_1 = 1$,
        $F_{n+2} = F_{n+1} + F_n$. Bewijs dat
        \[
        F^n = \begin{pmatrix} F_{n+1} & F_n\\ F_n &
        F_{n-1}\end{pmatrix} \quad (n \geq 1),
        \]
        leid de formule van Binet $F_n = \dfrac{\varphi^n -
        \psi^n}{\sqrt5}$ af, met $\varphi = \frac{1 + \sqrt5}2$ en
        $\psi = \frac{1 - \sqrt5}2$, en met vraag 7 de identiteit
        van Cassini $F_{n+1}F_{n-1} - F_n^2 = (-1)^n$.
\end{enumerate}

\medskip
\noindent\textbf{Deel III --- \omterm{pb:b1:matrices:1}{Lineaire recurrenties},
structureel.} Leg $s, p \in K$ met $p \neq 0$ vast, en zij $E_D$ de
\omterm{def:b1:logic:sets}{verzameling} van de rijen met $u_{n+2} = s\,u_{n+1} - p\,u_n$ voor
alle $n$.
\begin{enumerate}[resume]
  \item Toon aan dat $E_D$ een \omterm{def:b1:vspaces:def}{vectorruimte} van dimensie $2$ is
        (pas \cref{exo:b1:findim:10} aan).
  \item Toon aan dat de rij $(a_n)$ uit deel I het element van
        $E_D$ is met beginwaarden $0, 1$, en dat elke $u \in E_D$
        voldoet aan
        \[
        u_n = u_1\,a_n + u_0\,b_n \qquad (n \in \N),
        \]
        met $(b_n)$ als in deel I: de resten van de deling lossen
        \emph{alle} recurrenties tegelijk op.
  \item Zijn $\lambda \neq \mu$ de wortels van $D$, toon dan aan
        dat $\bigl((\lambda^n), (\mu^n)\bigr)$ een \omterm{def:b1:vspaces:free}{basis} van $E_D$
        is; is $D = (X-\lambda)^2$ met $\lambda \neq 0$, toon dan
        aan dat $\bigl((\lambda^n), (n\lambda^n)\bigr)$ er een is.
  \item Los volledig op: $u_{n+2} = u_{n+1} + 6u_n$, $u_0 = 1$,
        $u_1 = 8$; toets het antwoord aan $u_2$ en $u_3$.
  \item Zij $C = \begin{pmatrix} 0 & 1\\ -p & s\end{pmatrix}$ (de
        \emph{begeleidende matrix} van $D$). Toon aan dat
        \[
        \begin{pmatrix} u_{n}\\ u_{n+1}\end{pmatrix}
        = C^n \begin{pmatrix} u_0\\ u_1\end{pmatrix}
        \quad (u \in E_D),
        \]
        en dat $\operatorname{tr} C = s$ en $p(C) = p$: de
        recurrentie en de matrix dragen dezelfde \omterm{def:b1:poly:def}{veelterm} $D$.
\end{enumerate}

\medskip
\noindent\textbf{Deel IV --- Graad drie.} Zij $D_3 = X^3 - \alpha
X^2 - \beta X - \gamma$ en
\[
C_3 = \begin{pmatrix} 0 & 1 & 0\\ 0 & 0 & 1\\ \gamma & \beta &
\alpha \end{pmatrix}.
\]
\begin{enumerate}[resume]
  \item Toon aan dat $D_3(C_3) = 0$. \emph{(Bereken de \omterm{def:b1:linmaps:kerim}{beelden} van
        de canonieke basisvectoren onder machten van $C_3$: de
        \omterm{def:b1:logic:map}{afbeelding} van $C_3$ stuurt $e_1 \mapsto \dots \mapsto$ een
        combinatie die door de laatste rij wordt afgedwongen.)}
  \item Toon aan dat als $D_3$ drie verschillende wortels
        $\lambda_1, \lambda_2, \lambda_3$ heeft, de rest $R_n$ van
        $X^n$ gedeeld door $D_3$ de \emph{interpolant van Lagrange}
        is van de waarden $\lambda_i^n$ in de knooppunten
        $\lambda_i$ (\cref{thm:b1:poly:lagrange}); leid af dat elke
        ingang van $C_3^{\,n}$ een vaste \omterm{def:b1:linmaps:def}{lineaire} combinatie is van
        $\lambda_1^n, \lambda_2^n, \lambda_3^n$.
  \item Los op: $u_{n+3} = 2u_{n+2} + u_{n+1} - 2u_n$ met $u_0 =
        0$, $u_1 = 1$, $u_2 = 1$. \emph{(Ontbind $D_3 = (X - 1)(X +
        1)(X - 2)$.)} Controleer bij $u_3$.
  \item Bereken de rest van $X^n$ modulo $(X - \lambda)^3$
        \emph{(taylorontwikkeling van $X^n$ in $\lambda$)}, en leid
        een formule af voor $(\lambda I + N)^n$ wanneer $N^3 = 0$
        en $N$ met alles in zicht commuteert; toets haar aan het
        binomium.
  \item Toon aan dat voor $D_3$ met verschillende wortels de
        algemene oplossing van de recurrentie van orde $3$ gelijk
        is aan $u_n = c_1 \lambda_1^n + c_2\lambda_2^n +
        c_3\lambda_3^n$: bewijs dat de drie meetkundige rijen een
        \omterm{def:b1:vspaces:free}{basis} van de oplossingsruimte vormen. \emph{(Evalueer voor
        de vrijheid een nulcombinatie bij $n = 0, 1, 2$ en herken
        een interpolatiestelsel in de verschillende knooppunten
        $\lambda_i$.)}
\end{enumerate}

\medskip
\noindent\textbf{Deel V --- Fibonacci-dividenden, en synthese.}
\begin{enumerate}[resume]
  \item Bewijs dat $F_1 + F_2 + \dots + F_n = F_{n+2} - 1$.
  \item Leid uit $F^{m+n} = F^m F^n$ de somformule
        \[
        F_{m+n} = F_{m+1}F_n + F_m F_{n-1}
        \]
        af, en daaruit $F_{2n} = F_n(F_{n+1} + F_{n-1})$.
  \item Bewijs dat $F_n$ voor elke $n \geq 0$ het dichtstbijzijnde
        gehele getal bij $\varphi^n/\sqrt5$ is.
  \item Zij $t_n = \operatorname{tr}(F^n) = F_{n+1} + F_{n-1}$ (de
        \emph{lucasgetallen} $L_n$). Toon aan dat $t_{n+2} = t_{n+1}
        + t_n$, $t_1 = 1$, $t_2 = 3$, dat $L_n = \varphi^n +
        \psi^n$, en vind $F_{2n} = F_n L_n$ terug.
  \item Synthese, in vier zinnen: waarom de machten van een $2
        \times 2$ matrix in het vlak $\operatorname{Vect}(I, A)$
        van $\mathcal{M}_2(K)$ leven (welk dimensieargument een
        kwadratische identiteit garandeert, en welke expliciete
        identiteit deel II opleverde); hoe de euclidische deling
        machtsverheffen omzet in een recurrentie met twee termen;
        welke \omterm{def:b1:logic:statement}{uitspraak} van deze opgave het geval $n = 2$ is van
        een stelling die in alle dimensies geldt (benoem haar, en
        zeg waar zij in deze \omterm{def:b1:series:def}{reeks} wordt bewezen); en wat de
        constructie met de begeleidende matrix aan het \omterm{def:b1:linmaps:kerim}{beeld}
        toevoegt.
\end{enumerate}
\end{problem}
